\documentclass[11pt]{article}

\usepackage[a4paper,margin=1.05in]{geometry}
\usepackage{amsmath,amssymb,amsthm}
\usepackage{mathtools}
\usepackage[protrusion=true,expansion=false]{microtype}
\usepackage{enumitem}
\usepackage[dvipsnames]{xcolor}
\usepackage{tikz-cd}
\usepackage[hidelinks]{hyperref}

\hypersetup{
  colorlinks=true,
  linkcolor=BrickRed,
  citecolor=BrickRed,
  urlcolor=BrickRed
}

\theoremstyle{plain}
\newtheorem{theorem}{Theorem}
\newtheorem{proposition}[theorem]{Proposition}
\newtheorem{conjecture}[theorem]{Conjecture}
\theoremstyle{definition}
\newtheorem{definition}[theorem]{Definition}
\newtheorem{target}{Target}
\theoremstyle{remark}
\newtheorem{remark}[theorem]{Remark}

\newcommand{\Cost}{\mathsf{Cost}}
\newcommand{\ciGSLT}{\mathbf{ciGSLT}}
\newcommand{\ciGSLTtc}{\mathbf{ciGSLT}_{\mathrm{TC}}}
\newcommand{\HML}{\textsf{HML}}
\newcommand{\CP}{\textsf{CP}}
\newcommand{\quoteP}[1]{@#1}
\newcommand{\drop}[1]{{*}#1}
\newcommand{\bisim}{\;\approx\;}

\title{\textbf{Proof Search as Concurrent Reaction}\\[0.35em]
\large The Rho Calculus and the Category of Continued Interactive GSLTs\\
as a Substrate for Proof Objects with Internalized, Cost-Accounted Search}

\author{L.~G.~Meredith\\ \small F1R3FLY.io}
\date{\today}

\begin{document}
\maketitle

\begin{abstract}
Proof objects in contemporary proof assistants are variants of the typed
$\lambda$-calculus, so their dynamics --- normalization, equivalently cut
elimination --- is sequential and confluent. Concurrency, when these systems
are scaled out, is supplied by an orchestration layer built \emph{on top} in a
different model of computation, leaving proof search external to proof
dynamics. We argue that the category of continued interactive graph-structured
lambda theories ($\ciGSLT$) is a class of computational models better suited to
hosting proof objects, in which proof search is internalized as concurrent
proof dynamics, and that the rho calculus is the distinguished object of this
category on which every requirement is native. The base case of the proposal
already exists as a theorem --- the Caires--Pfenning--Wadler correspondence
between linear-logic proofs and session-typed processes, under which cut
elimination is communication --- and our contribution is two generalizations
stacked on that base: replacing the fixed logic with an arbitrary
$\ciGSLT$ whose propositions are modal (Hennessy--Milner) formulas over its
labelled transition system, and pushing from proof \emph{normalization} to
proof \emph{search}. The pivot is a single identification: an unproven
obligation is a pending communication. Reflection makes obligations
first-class, collapsing the object/meta distinction and placing the
orchestrator inside the dynamics. The cost-accounting endofunctor of the prior
notes then supplies search control in two registers --- an economic one in
which inference rules are priced and budgets replace coordinator algorithms,
and a biological one in which mortality (deadlock-by-starvation) prunes
unproductive branches without a coordinator and metabolism reallocates
proof-effort under selection. We close with three theorem-shaped targets:
encoding-adequacy of decentralized search, fairness-and-budget completeness in
the limit, and proof identity as bisimulation.
\end{abstract}

\section{The problem: orchestration outside the dynamics}

In Lean, Coq, Agda and their relatives, a proof object is a term of a typed
$\lambda$-calculus --- the Calculus of Inductive Constructions, or for some a
core in the spirit of System F --- and the Curry--Howard correspondence
identifies proof normalization with $\beta$-reduction and, equivalently, cut
elimination with reduction in the term calculus. This is the deepest and most
useful fact in the area, and it has a structural consequence that is rarely
named as a limitation because it is usually a virtue: proof \emph{dynamics is
sequential}. A proof term is a finished plan whose reduction is canonical up to
the Church--Rosser property; there is no room in the reduction relation itself
for genuine choice, and there are no open ends in a closed term that
\emph{interact}.

When such systems are scaled out --- parallel elaboration, distributed proof
obligations, hammer-style search across many backends, learned premise
selection feeding many provers --- concurrency is introduced by building an
orchestration system \emph{on top}, in a different model of computation: a
scheduler, a worklist, a message layer, a tactic interpreter that constructs
and rewrites terms from the outside. The orchestration lives beside the proof
dynamics, not within it. This is not an accident of engineering. It is forced
by the two properties just named, and the first contribution of this note is to
say precisely why.

We then argue the positive thesis. The category of continued interactive
graph-structured lambda theories, developed in the prior notes of this series
as the home of the cost-accounting endofunctor and its internalization
adjunction, is a class of models in which proof objects can be represented so
that proof search is \emph{internalized} as concurrent proof dynamics.
Interaction among proof objects yields a decentralized and distributed proof
search with no overarching coordinator. Because proof search is undecidable in
general, hosting proof representation and proof search in a category where cost
accounting is available --- and, on Turing-complete bases, \emph{internalizable}
--- supports both economically controlled search tactics and biologically and
evolutionarily inspired ones. The rho calculus is the distinguished object on
which all of this is native rather than aspirational.

\section{The base case is a theorem}

The proposal is not free-floating; its restriction to one fixed logic is an
established correspondence. Caires and Pfenning showed that a session type
discipline for the $\pi$-calculus corresponds exactly to the sequent calculus
of dual intuitionistic linear logic, with process reduction simulating proof
reduction and conversely~\cite{cairespfenning}. Wadler recast the classical case
as $\CP$, a calculus in which the propositions of classical linear logic
\emph{are} session types, communication \emph{is} cut elimination, and
deadlock freedom follows from the logical correspondence~\cite{wadler}. Both
descend from the older reading of linear logic proofs as
processes~\cite{girard}. So ``proof dynamics as concurrent interaction'' is,
at the level of normalization, an isomorphism rather than a hope.

Two features of this base case are exactly what our generalizations must lift.
First, it is pinned to a fixed logic: the propositions are linear-logic
formulas, and the calculus is a session-typed fragment of the
$\pi$-calculus~\cite{milnerparrowwalker}. Second, it concerns the dynamics of an
\emph{already-found} proof: cut elimination normalizes a proof that exists, it
does not search for one. The genuinely new step is the second, and most of the
technical weight of the programme falls there.

\section{Why the \texorpdfstring{$\lambda$}{lambda}-calculus forces the orchestrator outside}

It is worth stating the negative result crisply, since it is the entire
motivation. Two properties of the core calculus force proof search into a
separate metalanguage.

\emph{Confluence.} Reduction of a CIC term is canonical up to Church--Rosser,
so the dynamics has no room for genuine choice. But proof search \emph{is}
choice: which rule to apply, which hypothesis to use, which term to instantiate.
A confluent reduction relation cannot express the branching that search
explores, so the branching must be managed elsewhere.

\emph{Closedness.} A $\lambda$-term has no open ends that interact. A partial
proof can therefore be represented only as a term with a hole, and a
hole can be filled solely by an \emph{external} agent that inspects the term and
plugs it. There is no notion, internal to reduction, of an incomplete proof
reaching out for what it lacks.

Both missing affordances --- nondeterministic choice and interactive
open-endedness --- are precisely what a process calculus has and a
$\lambda$-calculus does not. Hence the meta/object split (Lean's elaboration
engine and metaprogramming monads, Coq's tactic language) is not a design taste
but the only place choice and open-endedness can live when proof objects are
confluent closed terms.

\section{Obligation is communication}

The pivot of the whole proposal is one identification.

\begin{quote}
\emph{An unproven obligation is a pending communication.}
\end{quote}

In the session-typed base case this is already literal: a free channel typed by
a formula $A$ \emph{is} an open hypothesis or goal of type $A$; the typing
context lists the obligations as channel types. Generalize the carrier from the
session-typed $\pi$-calculus to an arbitrary continued interactive GSLT and the
open goals of a partial proof become exactly the free names and pending inputs
of a proof-process. In the implementation vocabulary of the GSLT-to-rho
compiler of the prior notes, an open goal is a \emph{Hole}-bud and a rule scheme
awaiting instantiation is a \emph{Schema}-bud.

This identification dissolves the level distinction. Proof checking and proof
search are not two activities at two levels; they are \emph{one} dynamics
observed at different stages of saturation. Search is the dynamics of partial,
open proof-processes seeking interaction partners --- lemmas, axioms, rule
instances. Checking is the same dynamics after no obligations remain pending. A
completed proof is a configuration with no free obligations and no enabled
interaction: a normal form in the behavioral sense.

Reflection is what makes the identification \emph{usable} rather than merely
true. In the rho calculus, quotation $\quoteP{P}$ turns a process into a name and
dereference $\drop{x}$ turns a name back into the process it
quotes~\cite{meredithradestock}. A proof-process can therefore reify its own
open obligations as data, route them, inspect them, transform them, and run
them. This is what places the orchestrator \emph{inside} the dynamics: the
object language of proofs and the metalanguage of search-strategy live in one
calculus, because ``manipulate an obligation'' and ``manipulate a process as
data'' are the same operation. The session-typed base case did not need
reflection because normalization does not need it. Search does.

\section{Proof search as concurrent reaction}

The operational picture is the chemical-abstract-machine reading of
rewriting~\cite{berryboudol}, which is a hair's breadth from a GSLT presented as
types, equations, and rewrites~\cite{meseguer}. Goals, rules, lemmas and axioms
float as molecules in a solution and react when a redex is enabled. The
disjunctive structure of search --- the OR-node, the choice of which rule to try
--- is competing communication on a shared channel, the nondeterministic choice
the $\lambda$-calculus could not express. The conjunctive structure --- the
AND-node, the several premises a rule demands --- is parallel composition, and
in the rho calculus specifically it is a \emph{join pattern}:
\[
\textsf{for}\,(\,y_1 \leftarrow a_1\ \&\ \cdots\ \&\ y_n \leftarrow a_n\,)\,P
\]
fires only when all $n$ input obligations are simultaneously available. The
multi-premise rule is the join; the join firing is the coalition of subproofs
all succeeding together. This is the Curry--Howard reading between Coalition
Logic proofs and rho joins from the consensus-synthesis note, surfacing here as
the native realization of an AND-node: no coordinator assembles the conjunction,
the join does.

The reaction itself is the structural-congruence-plus-communication semantics of
the calculus: parallel composition with unit $0$ is the commutative monoid that
\emph{is} the solution, the communication rule is the reaction, and reflection
is what lets the molecules carry one another's recipes. The decentralized soup
is not a metaphor laid over the model; it is the operational semantics read
aloud, with no agenda, no global worklist, and locally enabled reactivity.

This regime has honest precedents, and the note should name them rather than
claim virgin ground. Concurrent constraint programming, with its ask/tell on a
shared store, is already a decentralized, coordinator-free proof search and a
foundational instance of ``proof search as the soup
dynamics''~\cite{saraswat}. The uniform-proofs and abstract-logic-programming
tradition makes proof search the operational semantics of a logic, with the
connectives read as fixed search instructions~\cite{miller}. The deltas of the
present proposal over both are exactly: reflection, which collapses object and
meta; and an internalized cost layer, which controls and selects.

\section{Rho versus CP: the inversion}

The reason the rho calculus beats Wadler's $\CP$ for \emph{this} task is close
to a clean inversion, and it is worth stating as such, because $\CP$'s strengths
are precisely the wrong strengths here.

$\CP$ is a triumph of correspondence at the level of normalization, and it pays
for that tightness in exactly the currency that ruins it as a search substrate.
Well-typed $\CP$ is deadlock-free by theorem; its cut-reduction fragment is
confluent and strongly normalizing; and what looks like choice is typed into
race-free internal and external choice~\cite{wadler,cairespfenning}. Those are
the properties that make cut elimination behave like a finished computation. But
proof \emph{search} requires the three things a type discipline of this kind is
built to prohibit:
\begin{enumerate}[leftmargin=2.2em,itemsep=0.25em]
  \item \emph{Nonconfluent choice} --- real races between candidate rules, since
        the order and outcome of choices is the content of search.
  \item \emph{Deadlock} --- because deadlock-by-starvation is the pruning
        mechanism; mortality is not a failure mode here, it is the garbage
        collector.
  \item \emph{Potential divergence} --- because the set of proofs is recursively
        enumerable but not recursive, so any sound and complete search must be
        permitted not to terminate.
\end{enumerate}
$\CP$ disciplines all three away statically, by typing, as prohibitions. The rho
calculus keeps all three live and disciplines them dynamically, by an economy.
For a search substrate this is the correct control surface: the pathologies on
which search depends are exactly the ones a type discipline exists to forbid.
The choice is not of a worse-behaved calculus but of the one whose control is in
the right place.

Reflection is the second decisive gap. $\CP$ and the intuitionistic session
calculus are first-order in the relevant sense: channels carry channels, not
proof terms one can run as data. To obtain proofs-as-data one bolts on a
higher-order layer or a metalanguage, reintroducing the object/meta split the
programme exists to dissolve. In the rho calculus, $\quoteP{(\cdot)}$ and
$\drop{(\cdot)}$ \emph{are} that layer, inside one calculus.

There is a third point, and it converts a liability into an asset. Moving proof
dynamics into a nonconfluent model costs the canonicity of normal forms, so
proof identity can no longer be $\beta\eta$-equality; the right notion becomes
bisimulation. In $\CP$ this is a genuine imposition, since the typed terms
\emph{want} an observational identity one would be overriding. In the rho
calculus, weak and barbed bisimulation is already the home equivalence; there is
no canonical-normal-form identity to surrender. So ``proof identity is
bisimulation'' stops being a redefinition and becomes the default reading, and
it lands directly on the modal logic: two proofs are the same exactly when no
Hennessy--Milner formula separates their behavior~\cite{hennessymilner}, which
is the equivalence the algorithmic-scientists thread already singles out as the
finest learnable one. The cost incurred on a generic $\ciGSLT$ is free on rho
specifically.

\section{The cost layer: economy and biology}

The cost-accounting endofunctor $\Cost$ of the prior notes adjoins a signature
monoid, a token stack, a signed-term wrapper, and replaces each rewrite by its
token-gated family, so that no step fires without consuming matching fuel; on
Turing-complete bases the internalization map $\Cost(G)\to G$ exhibits the
metered calculus as a behavioral retract of the bare one, up to weak
bisimulation. Applied to proof search this layer supplies control in two
registers.

\paragraph{Economic control.} Price each rule or tactic in phlogiston: a rewrite
is cheap, a decision-procedure call is dear. A goal-process carries a budget and
can fire only rules it can afford. Iterative deepening, best-first, and
breadth-versus-depth become \emph{budget policies} rather than coordinator
algorithms. Because cost is internalized, a goal-process can reason about its own
budget and reallocate --- self-budgeting search, a meta-level strategy with no
meta-level interpreter, since the meter lives in the same calculus as the proof.

\paragraph{Evolutionary control.} The mortal-computation note identified
\emph{death as deadlock-by-starvation}: a process whose token stack is exhausted
has no enabled rewrite and blocks permanently, and this configuration was native
to the un-metered model all along. Read as proof search, an unproductive branch
\emph{prunes itself} --- no coordinator decides to kill it; it fails to pay and
deadlocks. That is the decentralization claim made mechanical. \emph{Metabolism}
supplies the rest: a proven lemma is a high-energy metabolite, a stable
intermediate that downstream goal-processes consume; the splitters and joiners
that reconcile token granularity are catabolism and anabolism, reallocating
proof-fuel toward branches producing reusable structure; the gated rewrite is
the thermodynamic coupling forbidding work without payment. With reflection, a
tactic-process that closes goals cheaply can replicate into fresh contexts, with
the phlogiston economy as the selection pressure --- fitness internal, not an
externally imposed objective. The variation operator is then the
language-defining capability of the surrounding framework: mutating and
recombining the rewrite rules of a defined proof-language mutates the tactics,
with the language's types, equations, and rewrites playing the role of a genome
and the cost economy doing the selecting. This is Leo Buss's account of why
selection favours mortal organisms~\cite{buss}, transposed to proof search with
a precise substrate rather than an analogy.

\paragraph{One economy.} On the rho calculus, the blockchain gas economy, the
proof-search budget, and the metabolic energy are not three economies. They are
one phlogiston economy on one substrate. The node's gas \emph{is} a search
budget \emph{is} a metabolism. This is the payoff of doing the work on rho rather
than on a bespoke calculus: the control layer already exists, is deployed, and
is the same coin under all three readings.

\section{Soundness, proof identity, and completeness in the limit}

\paragraph{Soundness.} The search moves are the rewrites; the rewrites are the
bisimulation-preserving morphisms on which $\Cost$ was defined; preservation of
bisimulation is preservation of the generated Hennessy--Milner
property~\cite{hennessymilner}. Hence each search step is sound with respect to
the behavioral proposition it claims to preserve. The shift to accept is that
proofs here prove \emph{behavioral and modal} properties, not static type
inhabitation. This is the natural home for verifying concurrent and distributed
systems and a slightly awkward fit for, say, ordinary arithmetic; we lean into
the former.

\paragraph{Proof identity.} As noted, identity of proofs is bisimulation, which
on rho is the home equivalence rather than an imposition. Two proof-processes
denote the same proof when no modal formula separates them. This is coarser than
$\beta\eta$-equality and behavioral by design.

\paragraph{Completeness in the limit.} The cleanest statement uses the
internalization. Under fair scheduling and unbounded budget, the reachable
closed configurations of the bare calculus are exactly the recursively
enumerable set of proofs of the underlying theory --- semi-decidability realized
as reachability, the most a sound search can promise. Adding cost monotonically
\emph{restricts} reachability: every starvation-deadlock removes a branch, so
completeness degrades gracefully rather than catastrophically, and sending the
token supply to infinity recovers full reachability. The two knobs are
orthogonal and both must be maximized for completeness --- fairness rules out
\emph{scheduling} starvation, budget rules out \emph{resource} starvation --- and
the cost layer is exactly the dial between ``complete and possibly divergent''
and ``bounded and self-pruning.''

\begin{remark}
Cost accounting never converts an undecidable search into a decidable one. What
it converts is the \emph{character} of nontermination: from ``search may
silently diverge'' to ``search consumes a controllable budget at a controllable
rate.'' This is the anytime, resource-bounded reframing one actually wants in
practice, and it is the right place to host the undecidability of search.
\end{remark}

\section{Three theorem-shaped targets}

\begin{target}[Encoding and adequacy of decentralized search]
A proof-search strategy specified as a global traversal of an AND/OR space is
realized by a parallel composition of goal-processes and rule-processes in a
continued interactive GSLT, such that the reachable terminal configurations
correspond exactly to the proofs the strategy finds, and no process holds global
search state. The AND-nodes are join patterns, the OR-nodes are competing
communications, and the correspondence is up to bisimulation.
\end{target}

\begin{target}[Fairness-and-budget completeness]
Under a fairness condition on the labelled transition system and an unbounded
token budget, the concurrent search is complete with respect to the underlying
theory: every provable goal has a reachable closed configuration. Equivalently,
the reachable closed configurations coincide with the recursively enumerable set
of proofs. Finite budgets restrict this set monotonically, with each
starvation-deadlock excising one branch.
\end{target}

\begin{target}[Proof identity is bisimulation]
Two proof-processes denote the same proof if and only if they are
weakly bisimilar, equivalently if and only if no generated Hennessy--Milner
formula separates them. The characterization factors through the
modal logic so that ``same behavior'' and ``same proof'' provably coincide.
\end{target}

\section{The wider category and the agenda}

The reason ``rho sits inside $\ciGSLT$'' earns its keep is that the
cost endofunctor, the internalization adjunction, and the framework's ability to
\emph{define} a proof-language-specific GSLT and compile it down to rho all live
in that neighborhood. One can grow a logic-tailored GSLT, compile it into rho,
and run its proof search in the same soup as everything else; one can let the
language-defining layer evolve the rewrite rules, which is to say evolve the
tactics, with phlogiston as the selection pressure. The distinguished object
gives nativity --- concurrency and joins for AND/OR, reflection for the
meta-collapse, deadlock for mortality-as-pruning, bisimulation as the free
proof-identity, one economy for gas, budget, and metabolism --- and the
surrounding category gives room to explore: other objects, other logics, other
cost structures, with the bisimulation-preserving morphisms as the connecting
tissue.

The note that follows should open exactly here, on rho versus $\CP$ framed as
the inversion of \S6, establish the obligation-is-communication identification
of \S4 as its hinge, and only then generalize back up to the category, carrying
the three targets of \S10 as the formal obligations the programme takes on.

\begin{thebibliography}{99}

\bibitem{cairespfenning}
L.~Caires and F.~Pfenning.
\newblock Session types as intuitionistic linear propositions.
\newblock In \emph{CONCUR 2010}, LNCS \textbf{6269}, pp.~222--236. Springer, 2010.

\bibitem{wadler}
P.~Wadler.
\newblock Propositions as sessions.
\newblock \emph{Journal of Functional Programming} \textbf{24}(2--3):384--418,
2014. Conference version in \emph{ICFP 2012}, pp.~273--286.

\bibitem{girard}
J.-Y.~Girard.
\newblock Linear logic.
\newblock \emph{Theoretical Computer Science} \textbf{50}(1):1--101, 1987.

\bibitem{milnerparrowwalker}
R.~Milner, J.~Parrow, and D.~Walker.
\newblock A calculus of mobile processes, I.
\newblock \emph{Information and Computation} \textbf{100}(1):1--40, 1992.

\bibitem{meredithradestock}
L.~G.~Meredith and M.~Radestock.
\newblock A reflective higher-order calculus.
\newblock \emph{Electronic Notes in Theoretical Computer Science}
\textbf{141}(5):49--67, 2005.

\bibitem{hennessymilner}
M.~Hennessy and R.~Milner.
\newblock Algebraic laws for nondeterminism and concurrency.
\newblock \emph{Journal of the ACM} \textbf{32}(1):137--161, 1985.

\bibitem{miller}
D.~Miller, G.~Nadathur, F.~Pfenning, and A.~Scedrov.
\newblock Uniform proofs as a foundation for logic programming.
\newblock \emph{Annals of Pure and Applied Logic} \textbf{51}(1--2):125--157, 1991.

\bibitem{saraswat}
V.~A.~Saraswat, M.~Rinard, and P.~Panangaden.
\newblock Semantic foundations of concurrent constraint programming.
\newblock In \emph{POPL 1991}, pp.~333--352. ACM, 1991.

\bibitem{berryboudol}
G.~Berry and G.~Boudol.
\newblock The chemical abstract machine.
\newblock \emph{Theoretical Computer Science} \textbf{96}(1):217--248, 1992.

\bibitem{meseguer}
J.~Meseguer.
\newblock Conditional rewriting logic as a unified model of concurrency.
\newblock \emph{Theoretical Computer Science} \textbf{96}(1):73--155, 1992.

\bibitem{buss}
L.~W.~Buss.
\newblock \emph{The Evolution of Individuality}.
\newblock Princeton University Press, 1987.

\bibitem{series}
L.~G.~Meredith.
\newblock Cost accounting for the rho calculus; the cost endofunctor on
continued interactive GSLTs; mortal computation and metabolism.
\newblock F1R3FLY.io technical notes (this series), 2026.

\end{thebibliography}

\end{document}
